\documentclass[11pt]{article}
\usepackage[utf8]{inputenc}
\usepackage{amsmath,amssymb}
\usepackage{authblk}
\usepackage[margin=1in]{geometry}
\usepackage[hidelinks]{hyperref}
\usepackage{xurl}
\usepackage{setspace}

\title{Quantum Mechanics without Ontology?\\
\large Configuration Space, the Arena Problem, and the Road Not Taken in 1926}
\author[1]{Claes Johnson}
\affil[1]{KTH Royal Institute of Technology, SE-100 44 Stockholm, Sweden \\
\texttt{claesjohnson@gmail.com}}
\date{\today \\[0.4em]
\small\textit{Prepared in cooperation with Claude (Anthropic); the argument and claims are the author's, with the
AI assisting in drafting and in surveying the literature.}}

\begin{document}
\maketitle

\begin{abstract}
It is often said that quantum mechanics dispenses with ontology and that its predictive success licenses the
omission. This paper argues that the situation is both more precise and more serious. First, quantum mechanics
does not lack an ontology outright: under the preparation-independence assumption, the Pusey--Barrett--Rudolph
(PBR) theorem tells against a purely epistemic reading of the quantum state, so the state is naturally taken to
be ontic. The genuine difficulty is not the \emph{existence} of an ontology but its \emph{arena}: the state that
PBR renders real is a function on $3N$-dimensional configuration space, not on the three-dimensional space that
observers and instruments occupy. The interpretive troubles of quantum mechanics --- the measurement problem,
the status of collapse, the profusion of interpretations --- are, I argue, symptoms of this arena problem rather
than of a bare absence of the real. Second, I revisit the historical moment (spring 1926) at which the wave
function was lifted from three-dimensional space to configuration space, and the objections of Schr\"odinger and
Einstein were set aside; I ask why a commitment so basic was relinquished so rapidly. Third, I situate the
options in the contemporary debate between \emph{wave-function realism}, which reifies the configuration-space
object, and \emph{primitive ontology} approaches, which locate the fundamental beables in three-dimensional
space, and I present a recent three-dimensional charge-density formulation (RealQM) as one under-explored member
of the latter family --- a concrete realization of the road not taken. Finally, I use the case of quantum
computing to illustrate a practical corollary: technological programs can inherit the arena problem, staking
resources on the physical reality of the exponentially large amplitude, and the ontological status of that
amplitude is therefore not an idle question.
\end{abstract}

\setstretch{1.05}

\section{Introduction}\label{sec:intro}

A familiar defense of standard quantum mechanics holds that although the theory declines to say what is
physically real between measurements, its unrivaled record in predicting measurement statistics renders the
omission harmless. Two things are run together in this defense that ought to be kept apart: whether quantum
mechanics has an ontology at all, and, if it has one, whether that ontology is located where physics has always
located its subject matter --- in three-dimensional space. I shall argue that recent results and debates let us
be precise here. The quantum state is not plausibly a mere summary of an agent's knowledge (\S\ref{sec:pbr});
but the state it is natural to reify lives on configuration space, and it is this displacement of the fundamental
object from three-space to a high-dimensional arena --- not a bare absence of ontology --- that generates the
theory's foundational difficulties (\S\ref{sec:arena}). Read this way, the century-old interpretive impasse is an
\emph{arena problem}, and the space of responses is exactly the contemporary contest between wave-function
realism and primitive ontology (\S\ref{sec:realism}). I recall the historical juncture at which the arena was
chosen (\S\ref{sec:1926}), present a three-dimensional charge-density formulation as an under-explored primitive
ontology (\S\ref{sec:realqm}), and close with a practical corollary drawn from quantum computing
(\S\ref{sec:qc}).

\section{Ontology versus epistemology, and what PBR settles}\label{sec:pbr}

Two questions must be separated: \emph{ontology}, what exists; and \emph{epistemology}, what and how we may know
of it. On the epistemic reading associated with parts of the Copenhagen tradition, the quantum state represents
knowledge or information about outcomes rather than a physical property of an individual system. Harrigan and
Spekkens \cite{HarriganSpekkens} made this reading precise within an ontological-models framework, distinguishing
$\psi$-epistemic from $\psi$-ontic models. The Pusey--Barrett--Rudolph theorem \cite{PBR} then showed that, under
a preparation-independence assumption, $\psi$-epistemic models of the relevant kind cannot reproduce the quantum
predictions: distinct pure states cannot both correspond to the same underlying physical state. The natural
moral, widely (if not universally) drawn, is that the quantum state is \emph{ontic} --- a physical feature of the
system, not merely a catalogue of expectations.

This matters for the present argument in a specific way. It means that the slogan ``quantum mechanics has no
ontology'' is, taken literally, too strong: there is a real state, and it is the wave function. But PBR fixes
only \emph{that} the state is real, not \emph{where} it is real. The state it certifies as physical is
$\Psi(\mathbf x_1,\dots,\mathbf x_N)$, a field on $3N$-dimensional configuration space. The question that remains
--- and that PBR does not touch --- is whether an ontology carried by an object on configuration space can serve
the explanatory role physics has always demanded of its ontologies, which have without exception been ontologies
in three-dimensional space or four-dimensional spacetime. I turn to that question now.

\section{The arena problem}\label{sec:arena}

Every physical theory before quantum mechanics answered, first, the question of what exists, and located the
answer in ordinary space: masses at positions, fields on $\mathbb R^3$, molecules, a spacetime metric. The
predictive content followed as the computed behaviour of those spatially located things. The distinctive feature
of quantum mechanics is that its fundamental object is a function not on physical space but on a space whose
dimension grows with particle number. Kohn \cite{Kohn} put the working consequence bluntly in his Nobel lecture:
the many-electron wave function is not, for large $N$, a legitimate computational object, being unrecordable even
in principle beyond a modest number of particles. Whatever its ontic status, an object we cannot in principle
record and which does not inhabit the space of laboratories and observers sits awkwardly as the bearer of
physical reality.

The interpretive pathologies of quantum mechanics can be read as consequences of this displacement. The
\emph{measurement problem} arises because the link between the configuration-space state and definite
three-dimensional facts (pointer positions, tracks, clicks) is not supplied by the unitary dynamics and must be
added by hand, as collapse; the \emph{Heisenberg cut} is the unlocatable seam between the configuration-space
description and the three-dimensional description we actually read; the nonlocality that so troubled Einstein
\cite{EPR} is the difficulty of a correlated whole that has, in three-space, no separated parts to be local to.
On this diagnosis the trouble is not that there is nothing real, but that the something real is in the wrong
arena, and the recovery of ordinary three-dimensional facts from it is left unaccounted for.

\section{Wave-function realism and primitive ontology}\label{sec:realism}

The contemporary metaphysics of quantum mechanics is, in effect, an argument over the arena. \emph{Wave-function
realism} accepts the configuration-space object as fundamental and takes the high-dimensional space to be the real
arena of the world, our three-dimensional experience being derivative \cite{Albert,NeyAlbert,Ney,Wallace}. The
program is serious and has able defenders, but it purchases realism at the price of demoting three-dimensional
space to appearance and owes an account --- the ``arena'' or ``$3$D-emergence'' problem --- of how the manifest
spatial world is recovered. \emph{Primitive-ontology} approaches take the opposite route: they posit fundamental
beables in three-dimensional space or spacetime --- particles, a matter-density field, or flash events --- and
assign the wave function a nomological or guiding role rather than the role of the primary substance
\cite{Allori,DGZ,Maudlin,Bell}. On these views the wave function may still be ontic in PBR's sense while not
being the \emph{primitive} ontology; what is fundamentally real is in three dimensions.

The choice between these families is not decided by PBR, which speaks only to the reality of the state, not to its
arena. It is decided, if at all, by which program better discharges the explanatory duties of a physical theory:
recovering the three-dimensional phenomena, dissolving the measurement problem, and doing so without ad hoc
additions. It is against this backdrop that the historical contingency of the standard choice, and the
availability of three-dimensional alternatives, become salient.

\section{The choice of 1926}\label{sec:1926}

The arena was chosen quickly. In the first months of 1926 Schr\"odinger's wave mechanics represented the
one-electron atom by a field $\psi(\mathbf x)$ on ordinary space, and Schr\"odinger initially hoped to read
$|\psi|^2$ as a charge density there. The extension to $N$ interacting electrons required a single function
$\Psi(\mathbf x_1,\dots,\mathbf x_N)$ on configuration space; Born's statistical interpretation \cite{Born} then
read $|\Psi|^2$ as a probability over configurations. Schr\"odinger, in his correspondence with Lorentz
\cite{Przibram}, and Einstein, through the Solvay debates and later the EPR argument \cite{EPR}, both objected
that a theory whose fundamental object was a many-dimensional amplitude had relinquished a physical picture in
three-space; both were overruled as the G\"ottingen--Copenhagen synthesis consolidated
\cite{MehraRechenberg}. That the arena was fixed so rapidly, and the dissent of the theory's own principal
architects set aside, is historically striking and, I suggest, philosophically relevant: it shows the
configuration-space choice to have been a contingent turn rather than a forced one, taken under the pressure of
immediate predictive success and a positivist climate that treated the demand for an ontology as metaphysics in
the pejorative sense.

\section{A three-dimensional charge-density formulation}\label{sec:realqm}

That a three-dimensional alternative was available, and not merely desired, can be shown constructively. One may
represent the $N$ electrons of an atom or molecule not by a single amplitude on $3N$-space but by $N$
non-negative charge densities on ordinary space, supported on disjoint regions, and take the ground state to
minimize the ordinary Coulomb energy over these densities and their spatial partition; excited and
time-dependent behaviour follow by admitting complex phases. This ``RealQM'' construction \cite{RealQM} is a
member of the primitive-ontology family in the specific sense that its beables --- charge densities --- are
fields in three-dimensional space, and it is offered here not as an established replacement for standard quantum
mechanics but as an existence proof: a concrete, computable, three-dimensional ontology in the spirit
Schr\"odinger first sought. Two clarifications place it correctly in the debate. It is not a hidden-variable
theory: it adds no parameters beneath the amplitude (as de~Broglie--Bohm does) but replaces the amplitude, as
primary object, with charge in three-space. And it is compatible with PBR's ontic conclusion while declining
wave-function realism's arena: the real thing is a state, but a state of charge in three dimensions rather than a
field on configuration space. Its empirical reach and its open problems --- notably a full account of spin
exchange and of the statistics of measurements on entangled systems --- are documented elsewhere \cite{RealQM};
the philosophical point here is only that the primitive-ontology-in-3D option is live and was the road not taken.

\section{A practical corollary: quantum computing}\label{sec:qc}

The arena problem is not purely interpretive; a technological program can inherit it. The claimed advantage of
quantum computation rests on treating the joint state of $n$ qubits --- an amplitude in a $2^n$-dimensional
Hilbert space --- as a physical resource on which operations act ``in parallel.'' This is the
configuration-space object again, and whether the exponential advantage is physically real depends on whether
that object is a physical resource or a formal device for the statistics of a system that is, in three-space, of
modest size. The two horns of \S\ref{sec:realism} bear directly on the technology: on wave-function realism the
exponential resource is real and (decoherence permitting) exploitable; on a three-dimensional primitive ontology
there is no exponential thing in the device, and the scalable advantage lacks a substrate. I do not claim the
matter is settled --- small-scale devices demonstrably perform entangling operations and sampling tasks, and the
scalable, fault-tolerant regime is where the question lies. It is worth noting that skepticism about that regime
is not confined to ontological considerations: Kalai \cite{Kalai}, among others, argues on the basis of noise
that fault-tolerant quantum computation may face an in-principle barrier. The philosophical observation is
narrower and, I think, secure: the feasibility of scalable quantum computing is entangled with the ontological
status of the configuration-space state, so the question of arena that PBR leaves open is not academic but bears
on where a great deal of scientific effort is being directed.

\section{Conclusion}\label{sec:conclusion}

The claim that quantum mechanics abandons ontology should be replaced by a sharper one. The quantum state is
plausibly real (PBR); the difficulty is that it is real on configuration space, and the recovery of the
three-dimensional world from it is exactly what the theory does not supply --- the arena problem, of which the
measurement problem and its relatives are symptoms. The contemporary options are wave-function realism, which
keeps the arena and demotes three-space, and primitive ontology, which keeps three-space and demotes the wave
function to a non-primary role. The historical record shows the configuration-space arena to have been a rapid
and contested choice, not a necessity, and a constructive three-dimensional charge-density formulation shows a
primitive-ontology alternative to be available in fact and not only in aspiration. Whether one adopts it or not,
the lesson is that ``where is the quantum state real?'' is a substantive question with consequences reaching from
the measurement problem to the prospects of quantum technology, and that it was closed too early, too
confidently, in 1926.

\begin{thebibliography}{99}
\bibitem{HarriganSpekkens} N.~Harrigan and R.~W.~Spekkens, \emph{Einstein, incompleteness, and the epistemic view
of quantum states}, Found.\ Phys.\ \textbf{40}, 125 (2010).
\bibitem{PBR} M.~F.~Pusey, J.~Barrett, and T.~Rudolph, \emph{On the reality of the quantum state}, Nature
Physics \textbf{8}, 475 (2012).
\bibitem{Kohn} W.~Kohn, \emph{Nobel Lecture: Electronic structure of matter --- wave functions and density
functionals}, Rev.\ Mod.\ Phys.\ \textbf{71}, 1253 (1999).
\bibitem{EPR} A.~Einstein, B.~Podolsky, and N.~Rosen, \emph{Can quantum-mechanical description of physical
reality be considered complete?}, Phys.\ Rev.\ \textbf{47}, 777 (1935).
\bibitem{Born} M.~Born, \emph{Zur Quantenmechanik der Sto\ss vorg\"ange}, Z.\ Phys.\ \textbf{37}, 863 (1926).
\bibitem{Przibram} K.~Przibram (ed.), \emph{Letters on Wave Mechanics: Schr\"odinger, Planck, Einstein,
Lorentz}, Philosophical Library, 1967.
\bibitem{MehraRechenberg} J.~Mehra and H.~Rechenberg, \emph{The Historical Development of Quantum Theory},
Vol.~5 (Erwin Schr\"odinger and the Rise of Wave Mechanics), Springer, 1987.
\bibitem{Albert} D.~Z.~Albert, \emph{Elementary quantum metaphysics}, in J.~T.~Cushing, A.~Fine, and S.~Goldstein
(eds.), \emph{Bohmian Mechanics and Quantum Theory: An Appraisal}, Kluwer, 1996.
\bibitem{NeyAlbert} A.~Ney and D.~Z.~Albert (eds.), \emph{The Wave Function: Essays on the Metaphysics of Quantum
Mechanics}, Oxford University Press, 2013.
\bibitem{Ney} A.~Ney, \emph{The World in the Wave Function}, Oxford University Press, 2021.
\bibitem{Wallace} D.~Wallace, \emph{The Emergent Multiverse: Quantum Theory according to the Everett
Interpretation}, Oxford University Press, 2012.
\bibitem{Allori} V.~Allori, \emph{Primitive ontology and the structure of fundamental physical theories}, in
\cite{NeyAlbert}, pp.~58--75.
\bibitem{DGZ} D.~D\"urr, S.~Goldstein, and N.~Zangh\`i, \emph{Quantum equilibrium and the origin of absolute
uncertainty}, J.\ Stat.\ Phys.\ \textbf{67}, 843 (1992).
\bibitem{Maudlin} T.~Maudlin, \emph{Philosophy of Physics: Quantum Theory}, Princeton University Press, 2019.
\bibitem{Bell} J.~S.~Bell, \emph{Speakable and Unspeakable in Quantum Mechanics}, Cambridge University Press,
1987.
\bibitem{Kalai} G.~Kalai, \emph{The quantum computer puzzle}, Notices Amer.\ Math.\ Soc.\ \textbf{63}, 508
(2016).
\bibitem{RealQM} C.~Johnson, \emph{Quantum Mechanics as Multiphase 3D Continuum Mechanics}, and the RealQM
Gallery of simulations and companion articles, 2026; \url{https://claes542.github.io/RealMolecule/gallery.html}.
\end{thebibliography}

\end{document}
